\documentclass[aspectratio=169]{beamer}
\usetheme{Madrid}
\usecolortheme{default}

\usepackage{graphicx}
\usepackage{amsmath}
\usepackage{booktabs}
\usepackage{pgfpages}

% UGA Colors
\definecolor{UGARed}{HTML}{BA0C2F}
\definecolor{UGABlack}{HTML}{000000}

% Apply UGA colors to Madrid theme
\setbeamercolor{palette primary}{bg=UGARed,fg=white}
\setbeamercolor{palette secondary}{bg=UGABlack,fg=white}
\setbeamercolor{palette tertiary}{bg=UGARed,fg=white}
\setbeamercolor{title}{fg=UGARed}
\setbeamercolor{frametitle}{bg=UGARed,fg=white}
\setbeamercolor{structure}{fg=UGARed}
\setbeameroption{hide notes}

% UGA Logo on title page
\titlegraphic{\includegraphics[height=1.5cm]{assets/download.png}}
\logo{\includegraphics[height=1.5cm]{assets/download.png}}

% Title info
\title[WSC Continuation]{WSC Continuation: Lane-Level Calibration of Microscopic Freeway Simulation}
\subtitle{Optimization Systems and Scalable Calibration Pipeline}
\author{Korede Bishi}
\institute{Group Meeting Lab Series}
\date{February 2026}

\begin{document}

%==============================================================================
% SLIDE 1: Title
%==============================================================================
\begin{frame}
    \titlepage
    \note{Thanks for having me. Last time I presented the ANNSIM work, which was about validation—figuring out what matters when comparing simulation to real data. Today I'm picking up where that left off. ANNSIM answered: use lane-level metrics, use the fast integrator, and model arrivals carefully. Now the question is: given those answers, how do we actually tune the model? That's what this WSC continuation is about—optimization at scale.}
\end{frame}

%==============================================================================
% SLIDE 2: Study Corridor
%==============================================================================
\begin{frame}{Study Corridor: US-101 Freeway Segment}
    \begin{columns}[c]
        \begin{column}{0.65\textwidth}
            \centering
            \includegraphics[width=\textwidth]{assets/road-fig.pdf}
        \end{column}
        \begin{column}{0.35\textwidth}
            \small
            \begin{itemize}
                \item 4-lane freeway segment
                \item 5 mainline sensors (S1--S5)
                \item 2 on-ramps merging between sensors
                \item PeMS 15-min lane-level counts
                \item Total: $\sim$1.4 km
            \end{itemize}
        \end{column}
    \end{columns}
    \note{Before we get into the methods, let me ground you in the physical system. This is a 1.4 km stretch of US-101—four lanes, five sensor stations, two on-ramps. Each sensor is basically a cross-section where we count how many vehicles pass per lane every 15 minutes. The ramps are where things get interesting—new vehicles merge in, and that's where lane distributions start to shift. This is real PeMS data, not synthetic.}
\end{frame}

%==============================================================================
% SLIDE 3: ANNSIM Results Bridge
%==============================================================================
\begin{frame}{ANNSIM Takeaway: Lane-Level Validation}
    \begin{columns}[c]
        \begin{column}{0.5\textwidth}
            \textbf{What ANNSIM established:}
            \begin{itemize}
                \item Integrator choice: Ballistic is sufficient (3$\times$ faster, $<$1\% accuracy loss)
                \item Arrival process: Shifted Erlang-2 reduces error by 28\%
                \item Validation level: Lane-level metrics expose what corridor averages hide
            \end{itemize}

            \vspace{0.8em}
            \textbf{What ANNSIM left open:}
            \begin{itemize}
                \item We used default IDM parameters
                \item No systematic calibration yet
                \item Lane-level distribution shift remains
            \end{itemize}
        \end{column}
        \begin{column}{0.5\textwidth}
            \centering
            \includegraphics[width=\textwidth]{assets/fig3_micro_flow_validation.pdf}
        \end{column}
    \end{columns}
    \note{So here's where we came from. ANNSIM was about validation—what do you actually compare, and does it matter how you compute the physics? Turns out integrator choice barely matters, but the arrival process does. And critically, when you look lane by lane instead of corridor totals, you see patterns that aggregate metrics completely miss. Fast lanes behave differently from slow lanes near merges. That's valuable, but we didn't calibrate anything—we just used textbook IDM values. So the simulation works, but it's not tuned. That's the gap this work fills.}
\end{frame}

%==============================================================================
% SLIDE 4: Motivation
%==============================================================================
\begin{frame}{The Calibration Challenge}
    \textbf{Core challenge:}
    \begin{itemize}
        \item Expensive black-box simulation objective
        \item Noisy, constrained, nonconvex parameter search
        \item Many repeated runs needed for fair benchmarking
        \item Lane-level distribution shift is the failure mode hidden by aggregate validation
    \end{itemize}

    \vspace{0.7em}
    \textbf{ANNSIM gave us two computational priors:}
    \begin{itemize}
        \item Ballistic integrator is runtime-efficient and sufficiently accurate
        \item Lane-level metrics are the right objective target
    \end{itemize}
    \note{So what makes this hard? The simulation is expensive—each run takes minutes. The objective function is noisy because of stochastic arrivals. The parameter space is constrained and nonconvex, so gradient descent doesn't just work out of the box. And here's the kicker: if you only look at corridor totals, you might think you're doing great, but individual lanes can be way off. That's the distribution shift problem. ANNSIM told us lane-level is the right target. Now we need to actually optimize for it.}
\end{frame}

%==============================================================================
% SLIDE 5: Research Question
%==============================================================================
\begin{frame}{WSC Continuation Question}
    \centering
    \large
    \textbf{Which model + optimizer pipeline best reduces} \\
    \textbf{lane-level distribution shift under realistic compute budget?}

    \vspace{1em}
    \normalsize
    \begin{itemize}
        \item Models: \textbf{IDM, Gipps, Krauss}
        \item Optimizers: \textbf{SPSA, SPSA-Mo, GA, Nelder-Mead}
        \item Validation target: lane-level flow and speed metrics
    \end{itemize}
    \note{This is the central question for WSC. We have three car-following models—IDM, Gipps, Krauss—and four optimizers with different search strategies. The goal isn't just to get low error. It's to reduce the gap between what the simulation says happens per lane and what the real sensors observed. That's distribution shift. We want to find which combination does that best, under a realistic compute budget.}
\end{frame}

%==============================================================================
% SLIDE 6: Benchmark Space
%==============================================================================
\begin{frame}{Benchmark Design Space}
    \begin{tabular}{l p{0.72\textwidth}}
        \toprule
        \textbf{Dimension} & \textbf{Definition} \\
        \midrule
        Models & IDM, Gipps, Krauss \\
        Optimizers & SPSA, SPSA-Mo, GA, Nelder-Mead \\
        Objective & Fitness \(= 0.5 \times\) Flow NRMSE \(+ 0.5 \times\) Speed NRMSE \\
        Constraints & Shared bounds and common initialization policy \\
        Shift diagnostics & Per-lane error profiles and cross-lane mismatch trends \\
        Output & Accuracy, distribution-shift diagnostics, convergence, runtime \\
        \bottomrule
    \end{tabular}

    \vspace{0.7em}
    \textbf{Fixed controls:}
    \begin{itemize}
        \item Ballistic integrator for production benchmarking
        \item Same data window and lane-level evaluation protocol
    \end{itemize}
    \note{This table lays out the full benchmark space. Three models, four optimizers, a combined flow-and-speed fitness function. The key constraint is fairness—same bounds, same starting conditions, same data window. That way when we compare results, we're comparing apples to apples. And critically, we track per-lane error profiles, not just aggregate fitness. That's how we catch distribution shift.}
\end{frame}

%==============================================================================
% SLIDE 7: Optimizers
%==============================================================================
\begin{frame}{Optimization Algorithms Under Test}
    \begin{tabular}{l p{0.7\textwidth}}
        \toprule
        \textbf{Optimizer} & \textbf{Why included} \\
        \midrule
        SPSA & Efficient stochastic gradient approximation for noisy simulation objectives \\
        SPSA-Mo & Momentum variant to improve stability and convergence speed \\
        GA & Population-based global search over nonconvex parameter space \\
        Nelder-Mead & Gradient-free simplex method as a strong deterministic baseline \\
        \bottomrule
    \end{tabular}
    \note{Why these four? They represent different search philosophies. SPSA approximates gradients with random perturbations—cheap and effective for noisy objectives. SPSA-Mo adds momentum to stabilize convergence. GA is population-based, so it explores globally but costs more evaluations. Nelder-Mead is a classic simplex method—no gradients, no randomness, just geometric search. Each has different strengths, and we want to know which works best for this problem.}
\end{frame}

%==============================================================================
% SLIDE 8: Runtime Engineering
%==============================================================================
\begin{frame}{Runtime Engineering}
    \textbf{Single 12-hour simulation runtime:}
    \begin{itemize}
        \item Ballistic: \(\sim\)15 minutes (previous baseline)
        \item Dormand-Prince: \(\sim\)55 minutes
    \end{itemize}

    \vspace{0.7em}
    \textbf{After virtual-thread + actor/director refactor:}
    \begin{itemize}
        \item Ballistic reduced to \(\sim\)2.5 minutes per 12-hour run
        \item ~6x speedup versus prior Ballistic baseline
    \end{itemize}

    \vspace{0.7em}
    \textbf{Impact:}
    \begin{itemize}
        \item Large optimizer\(\times\)model batches become practical on HPC job arrays
        \item Enables repeated runs for stability and reproducibility checks
    \end{itemize}
    \note{This is where engineering meets research. A 12-hour simulation used to take 55 minutes with Dormand-Prince. We moved to Ballistic, which got us to 15 minutes. Then we refactored the runtime—virtual threads, actor-director pattern—and now we're at 2.5 minutes. That's a 6x speedup over our previous baseline. Why does this matter? Because optimization means hundreds of simulation runs. If each run takes an hour, you can't iterate. At 2.5 minutes, you can actually explore the parameter space.}
\end{frame}

%==============================================================================
% SLIDE 9: Objective and Protocol
%==============================================================================
\begin{frame}{Objective Function and Fairness Protocol}
    \textbf{Decision variables (example for IDM):}
    \[
    \theta = [s_0,\; a_{\max},\; b,\; T,\; \tau]
    \]

    \textbf{Objective (minimize):}
    \[
    \text{Fitness} = 0.5 \times \text{Flow NRMSE} + 0.5 \times \text{Speed NRMSE}
    \]

    \textbf{Protocol for fair comparison:}
    \begin{itemize}
        \item Shared bounds and starting conditions for fair optimizer comparison
        \item Same simulation horizon and detector set across methods
        \item Lane-level metrics reported per sensor and per lane
        \item Same reporting format for convergence and final fit
        \item Track per-lane residual patterns, not just aggregate fitness
    \end{itemize}
    \note{Here's the nuts and bolts. For IDM, we're tuning five parameters—minimum spacing, max acceleration, comfortable deceleration, desired time headway, and reaction time. The fitness function combines flow error and speed error, equally weighted. Quick aside on why we use NRMSE: flow is in vehicles—maybe 200 per interval. Speed is in km/h—maybe 80. If I say flow error is 15 and speed error is 5, which is worse? You can't compare them directly—it's like asking whether 15 apples is more than 5 oranges. NRMSE turns both into percentages by dividing by the observed mean. Now 15 out of 200 is 7.5 percent, and 5 out of 80 is 6.25 percent. Suddenly you can add them, weight them, and optimize a single number. That's the whole point—making different measurements speak the same language. Now, the real protocol is about fairness: same bounds, same seeds, same data window across all experiments. And we don't just report aggregate fitness—we break it down by lane. That's how you catch hidden failures. A model might look good overall but completely miss Lane 4 near the merge.}
\end{frame}

%==============================================================================
% SLIDE 10: Experiment Matrix
%==============================================================================
\begin{frame}{Experiment Matrix}
    \centering
    \begin{tabular}{lcccc}
        \toprule
        \textbf{Model} & \textbf{SPSA} & \textbf{SPSA-Mo} & \textbf{GA} & \textbf{Nelder-Mead} \\
        \midrule
        IDM & \(\checkmark\) & \(\checkmark\) & \(\checkmark\) & \(\checkmark\) \\
        Gipps & \(\checkmark\) & \(\checkmark\) & \(\checkmark\) & \(\checkmark\) \\
        Krauss & \(\checkmark\) & \(\checkmark\) & \(\checkmark\) & \(\checkmark\) \\
        \bottomrule
    \end{tabular}

    \vspace{1em}
    \textbf{Total core combinations:} \(3 \times 4 = 12\)

    \vspace{0.4em}
    \textit{Optional repeats by random seed for stability analysis}
    \note{So this is the full matrix—12 core combinations. Three models times four optimizers. Each cell is a complete calibration run. We'll also do repeated runs with different random seeds to check stability. Some optimizers are stochastic, so one run isn't enough to know if the result is reliable or just lucky.}
\end{frame}

%==============================================================================
% SLIDE 11: Current Status
%==============================================================================
\begin{frame}{Current Status and Immediate Next Steps}
    \textbf{Confirmed current status:}
    \begin{itemize}
        \item ANNSIM baseline and integrator decision are complete
        \item Calibration framework and optimizer runners are in place
        \item Runtime pipeline upgraded with virtual threads and actor/director refactor
    \end{itemize}

    \vspace{0.7em}
    \textbf{Important note:}
    \begin{itemize}
        \item Full WSC optimization benchmark results are \textbf{not yet available}
        \item This phase is currently in active run and analysis
    \end{itemize}

    \vspace{0.4em}
    \textbf{Immediate next steps:}
    \begin{itemize}
        \item Execute full benchmark job array on Sapelo2
        \item Produce lane-level comparison tables and convergence plots
        \item Add lane-level distribution-shift summary plots to benchmark reports
        \item Report best model+optimizer combination in next update
    \end{itemize}
    \note{Let me be clear about where we are. The infrastructure is done—framework, optimizers, runtime pipeline. But the full benchmark hasn't finished running yet. I'm not going to show you fabricated numbers. What I can tell you is: the setup is solid, and results are coming. By the next update, we'll have the full comparison across all 12 combinations.}
\end{frame}

%==============================================================================
% SLIDE 12: Planned WSC Deliverables
%==============================================================================
\begin{frame}{Planned WSC Deliverables From This Phase}
    \begin{enumerate}
        \item Best model+optimizer recommendation for lane-level fidelity
        \item Comparative table of accuracy and runtime across the benchmark matrix
        \item Convergence behavior analysis (speed, stability, consistency)
        \item Lane-level distribution-shift analysis and mitigation summary
        \item Reproducibility checklist: bounds, seeds, budgets, and run metadata
    \end{enumerate}

    \vspace{0.8em}
    \textbf{Success criterion:} reproducible model+optimizer comparison with transparent lane-level metrics.
    \note{This is what the final deliverable looks like. Not just "which optimizer wins" but a full breakdown: accuracy, runtime, convergence behavior, and—critically—lane-level distribution shift. The goal is reproducibility. Someone else should be able to take our bounds, our seeds, our protocol, and get the same results. That's the bar we're aiming for.}
\end{frame}

%==============================================================================
% SLIDE 13: Closing
%==============================================================================
\begin{frame}{Closing}
    \textbf{Summary:}
    \begin{itemize}
        \item ANNSIM established the validation and integrator baseline
        \item WSC continuation focuses on optimization, scalability, and lane-level distribution-shift reduction
        \item Next report will present full comparative results
    \end{itemize}

    \vspace{0.8em}
    \textbf{Core contribution:} A scalable calibration pipeline that exposes and reduces hidden lane-level distribution failures in traffic simulation.

    \vspace{0.8em}
    \centering
    \Large\textbf{Questions?}
    \note{So to wrap up: ANNSIM gave us the validation foundation—what to measure and how. This WSC phase is about optimization—finding the best model-optimizer pairing under realistic compute constraints. The core contribution is a pipeline that doesn't just minimize aggregate error, but exposes and reduces lane-level distribution shift. That's the gap between what the simulation predicts per lane and what actually happens. If someone asks about SUMO: I'm not replacing it. I'm building a methodology for calibration and validation that makes simulation evidence more trustworthy for downstream applications like AV testing.}
\end{frame}

\end{document}
